\section{Distinguished triangles in the homotopy category}
\label{section-homotopy-triangulated}

\noindent
Since we want our boundary maps in long exact sequences of cohomology
to be given by the maps in the snake lemma without signs we define
distinguished triangles in the homotopy category as follows.

\begin{definition}
\label{definition-distinguished-triangle}
Let $\mathcal{A}$ be an additive category.
A triangle $(X, Y, Z, f, g, h)$ of $K(\mathcal{A})$ is
called a {\it distinguished triangle of $K(\mathcal{A})$}
if it is isomorphic to the triangle associated to
a termwise split exact sequence of complexes, see Definition
\ref{definition-split-ses}.
Same definition for $K^{+}(\mathcal{A})$, $K^{-}(\mathcal{A})$, and
$K^b(\mathcal{A})$.
\end{definition}

\noindent
Note that according to Lemma \ref{lemma-the-same-up-to-isomorphisms}
a triangle of the form $(K^\bullet, L^\bullet, C(f)^\bullet, f, i, -p)$
is a distinguished triangle.
This does indeed lead to a triangulated category, see
Proposition \ref{proposition-homotopy-category-triangulated}.
Before we can prove the proposition we need one more lemma
in order to be able to prove TR4.

\begin{lemma}
\label{lemma-two-split-injections}
Let $\mathcal{A}$ be an additive category. Suppose that
$\alpha : A^\bullet \to B^\bullet$ and $\beta : B^\bullet \to C^\bullet$
are split injections of complexes. Then there exist distinguished triangles
$(A^\bullet, B^\bullet, Q_1^\bullet, \alpha, p_1, d_1)$,
$(A^\bullet, C^\bullet, Q_2^\bullet, \beta \circ \alpha, p_2, d_2)$
and
$(B^\bullet, C^\bullet, Q_3^\bullet, \beta, p_3, d_3)$
for which TR4 holds.
\end{lemma}

\begin{proof}
Say $\pi_1^n : B^n \to A^n$, and $\pi_3^n : C^n \to B^n$ are the splittings.
Then also $A^\bullet \to C^\bullet$ is a split injection with splittings
$\pi_2^n = \pi_1^n \circ \pi_3^n$. Let us write $Q_1^\bullet$, $Q_2^\bullet$
and $Q_3^\bullet$ for the ``quotient'' complexes. In other words,
$Q_1^n = \Ker(\pi_1^n)$, $Q_3^n = \Ker(\pi_3^n)$ and
$Q_2^n = \Ker(\pi_2^n)$. Note that the kernels exist. Then
$B^n = A^n \oplus Q_1^n$ and $C^n = B^n \oplus Q_3^n$, where we think of $A^n$
as a subobject of $B^n$ and so on. This implies
$C^n = A^n \oplus Q_1^n \oplus Q_3^n$. Note that
$\pi_2^n = \pi_1^n \circ \pi_3^n$ is zero on both $Q_1^n$ and $Q_3^n$. Hence
$Q_2^n = Q_1^n \oplus Q_3^n$. Consider the commutative diagram
$$
\begin{matrix}
0 & \to & A^\bullet & \to & B^\bullet & \to & Q_1^\bullet & \to & 0 \\
  &     & \downarrow &     & \downarrow &     & \downarrow  & \\
0 & \to & A^\bullet & \to & C^\bullet & \to & Q_2^\bullet & \to & 0 \\
  &     & \downarrow &     & \downarrow &     & \downarrow  & \\
0 & \to & B^\bullet & \to & C^\bullet & \to & Q_3^\bullet & \to & 0
\end{matrix}
$$
The rows of this diagram are termwise split exact sequences, and
hence determine distinguished triangles by
definition. Moreover downward arrows in the diagram above
are compatible with the chosen splittings and hence
define morphisms of triangles
$$
(A^\bullet \to B^\bullet \to Q_1^\bullet \to A^\bullet[1])
\longrightarrow
(A^\bullet \to C^\bullet \to Q_2^\bullet \to A^\bullet[1])
$$
and
$$
(A^\bullet \to C^\bullet \to Q_2^\bullet \to A^\bullet[1])
\longrightarrow
(B^\bullet \to C^\bullet \to Q_3^\bullet \to B^\bullet[1]).
$$
Note that the splittings $Q_3^n \to C^n$
of the bottom split sequence in the diagram provides a splitting
for the split sequence
$0 \to Q_1^\bullet \to Q_2^\bullet \to Q_3^\bullet \to 0$
upon composing with $C^n \to Q_2^n$. It follows easily from this
that the morphism $\delta : Q_3^\bullet \to Q_1^\bullet[1]$
in the corresponding distinguished triangle
$$
(Q_1^\bullet \to Q_2^\bullet \to Q_3^\bullet \to Q_1^\bullet[1])
$$
is equal to the composition $Q_3^\bullet \to B^\bullet[1] \to Q_1^\bullet[1]$.
Hence we get a structure as in the conclusion of axiom TR4.
\end{proof}

\begin{proposition}
\label{proposition-homotopy-category-triangulated}
Let $\mathcal{A}$ be an additive category.
The category $K(\mathcal{A})$ of complexes up to
homotopy with its natural translation functors
and distinguished triangles as defined above
is a triangulated category.
\end{proposition}

\begin{proof}
Proof of TR1. By definition every triangle isomorphic to a distinguished
one is distinguished. Also, any triangle $(A^\bullet, A^\bullet, 0, 1, 0, 0)$
is distinguished since $0 \to A^\bullet \to A^\bullet \to 0 \to 0$ is
a termwise split sequence of complexes. Finally, given any morphism of
complexes $f : K^\bullet \to L^\bullet$ the triangle
$(K, L, C(f), f, i, -p)$ is distinguished by
Lemma \ref{lemma-the-same-up-to-isomorphisms}.

\medskip\noindent
Proof of TR2. Let $(X, Y, Z, f, g, h)$ be a triangle.
Assume $(Y, Z, X[1], g, h, -f[1])$ is distinguished.
Then there exists a termwise split sequence of complexes
$A^\bullet \to B^\bullet \to C^\bullet$ such that the associated
triangle $(A^\bullet, B^\bullet, C^\bullet, \alpha, \beta, \delta)$
is isomorphic to $(Y, Z, X[1], g, h, -f[1])$. Rotating back we see
that $(X, Y, Z, f, g, h)$ is isomorphic to
$(C^\bullet[-1], A^\bullet, B^\bullet, -\delta[-1], \alpha, \beta)$.
It follows from Lemma \ref{lemma-rotate-triangle} that the triangle
$(C^\bullet[-1], A^\bullet, B^\bullet, \delta[-1], \alpha, \beta)$
is isomorphic to
$(C^\bullet[-1], A^\bullet, C(\delta[-1])^\bullet, \delta[-1], i, p)$.
Precomposing the previous isomorphism of triangles with $-1$ on $Y$
it follows that $(X, Y, Z, f, g, h)$ is isomorphic to
$(C^\bullet[-1], A^\bullet, C(\delta[-1])^\bullet, \delta[-1], i, -p)$.
Hence it is distinguished by
Lemma \ref{lemma-the-same-up-to-isomorphisms}.
On the other hand, suppose that $(X, Y, Z, f, g, h)$ is distinguished.
By Lemma \ref{lemma-the-same-up-to-isomorphisms} this means that it is
isomorphic to a triangle of the form
$(K^\bullet, L^\bullet, C(f), f, i, -p)$ for some morphism of
complexes $f$. Then the rotated triangle $(Y, Z, X[1], g, h, -f[1])$ is
isomorphic to $(L^\bullet, C(f), K^\bullet[1], i, -p, -f[1])$ which is
isomorphic to the triangle $(L^\bullet, C(f), K^\bullet[1], i, p, f[1])$.
By Lemma \ref{lemma-rotate-cone} this triangle is distinguished.
Hence $(Y, Z, X[1], g, h, -f[1])$ is distinguished as desired.

\medskip\noindent
Proof of TR3. Let
$(X, Y, Z, f, g, h)$ and $(X', Y', Z', f', g', h')$
be distinguished triangles of $K(\mathcal{A})$
and let $a : X \to X'$ and $b : Y \to Y'$ be morphisms
such that $f' \circ a = b \circ f$. By
Lemma \ref{lemma-the-same-up-to-isomorphisms} we may assume that
$(X, Y, Z, f, g, h) = (X, Y, C(f), f, i, -p)$ and
$(X', Y', Z', f', g', h') = (X', Y', C(f'), f', i', -p')$.
At this point we simply apply Lemma \ref{lemma-functorial-cone}
to the commutative diagram given by $f, f', a, b$.

\medskip\noindent
Proof of TR4. At this point we know that $K(\mathcal{A})$
is a pre-triangulated category. Hence we can use
Lemma \ref{lemma-easier-axiom-four}. Let $A^\bullet \to B^\bullet$
and $B^\bullet \to C^\bullet$ be composable morphisms of
$K(\mathcal{A})$. By Lemma \ref{lemma-sequence-maps-split} we may assume that
$A^\bullet \to B^\bullet$ and $B^\bullet \to C^\bullet$
are split injective morphisms. In this case the result follows
from Lemma \ref{lemma-two-split-injections}.
\end{proof}

\begin{remark}
\label{remark-boundedness-conditions-triangulated}
Let $\mathcal{A}$ be an additive category.
Exactly the same proof as the proof of
Proposition \ref{proposition-homotopy-category-triangulated}
shows that the categories
$K^{+}(\mathcal{A})$, $K^{-}(\mathcal{A})$, and $K^b(\mathcal{A})$
are triangulated categories. Namely, the cone of a morphism between
bounded (above, below) is bounded (above, below).
But we prove below that these are triangulated subcategories
of $K(\mathcal{A})$ which gives another proof.
\end{remark}

\begin{lemma}
\label{lemma-bounded-triangulated-subcategories}
Let $\mathcal{A}$ be an additive category. The categories
$K^{+}(\mathcal{A})$, $K^{-}(\mathcal{A})$, and $K^b(\mathcal{A})$
are full triangulated subcategories of $K(\mathcal{A})$.
\end{lemma}

\begin{proof}
Each of the categories mentioned is a full additive subcategory.
We use the criterion of
Lemma \ref{lemma-triangulated-subcategory}
to show that they are triangulated subcategories.
It is clear that each of the categories
$K^{+}(\mathcal{A})$, $K^{-}(\mathcal{A})$, and $K^b(\mathcal{A})$
is preserved under the shift functors $[1], [-1]$.
Finally, suppose that $f : A^\bullet \to B^\bullet$ is a morphism
in $K^{+}(\mathcal{A})$, $K^{-}(\mathcal{A})$, or $K^b(\mathcal{A})$.
Then $(A^\bullet, B^\bullet, C(f)^\bullet, f, i, -p)$ is a distinguished
triangle of $K(\mathcal{A})$ with $C(f)^\bullet \in K^{+}(\mathcal{A})$,
$K^{-}(\mathcal{A})$, or $K^b(\mathcal{A})$ as is clear from the construction
of the cone. Thus the lemma is proved. (Alternatively,
$K^\bullet \to L^\bullet$ is isomorphic to an termwise split injection
of complexes in $K^{+}(\mathcal{A})$, $K^{-}(\mathcal{A})$, or
$K^b(\mathcal{A})$, see
Lemma \ref{lemma-make-injective}
and then one can directly take the associated
distinguished triangle.)
\end{proof}

\begin{lemma}
\label{lemma-additive-exact-homotopy-category}
Let $\mathcal{A}$, $\mathcal{B}$ be additive categories.
Let $F : \mathcal{A} \to \mathcal{B}$ be an additive functor.
The induced functors
$$
\begin{matrix}
F : K(\mathcal{A}) \longrightarrow K(\mathcal{B}) \\
F : K^{+}(\mathcal{A}) \longrightarrow K^{+}(\mathcal{B}) \\
F : K^{-}(\mathcal{A}) \longrightarrow K^{-}(\mathcal{B}) \\
F : K^b(\mathcal{A}) \longrightarrow K^b(\mathcal{B})
\end{matrix}
$$
are exact functors of triangulated categories.
\end{lemma}

\begin{proof}
Suppose $A^\bullet \to B^\bullet \to C^\bullet$
is a termwise split sequence of complexes of $\mathcal{A}$ with splittings
$(s^n, \pi^n)$ and associated morphism $\delta : C^\bullet \to A^\bullet[1]$,
see Definition \ref{definition-split-ses}. Then
$F(A^\bullet) \to F(B^\bullet) \to F(C^\bullet)$
is a termwise split sequence of complexes with splittings
$(F(s^n), F(\pi^n))$ and associated morphism
$F(\delta) : F(C^\bullet) \to F(A^\bullet)[1]$.
Thus $F$ transforms distinguished triangles into distinguished triangles.
\end{proof}

\begin{lemma}
\label{lemma-improve-distinguished-triangle-homotopy}
Let $\mathcal{A}$ be an additive category. Let
$(A^\bullet, B^\bullet, C^\bullet, a, b, c)$ be a distinguished triangle in
$K(\mathcal{A})$. Then there exists an isomorphic distinguished triangle
$(A^\bullet, (B')^\bullet, C^\bullet, a', b', c)$ such that
$0 \to A^n \to (B')^n \to C^n \to 0$ is a split short exact sequence
for all $n$.
\end{lemma}

\begin{proof}
We will use that $K(\mathcal{A})$ is a triangulated category by
Proposition \ref{proposition-homotopy-category-triangulated}.
Let $W^\bullet$ be the cone on $c : C^\bullet \to A^\bullet[1]$ with its maps
$i : A^\bullet[1] \to W^\bullet$ and $p : W^\bullet \to C^\bullet[1]$. Then
$(C^\bullet, A^\bullet[1], W^\bullet, c, i, -p)$ is a distinguished triangle
by Lemma \ref{lemma-the-same-up-to-isomorphisms}. Rotating backwards twice
we see that $(A^\bullet, W^\bullet[-1], C^\bullet, -i[-1], p[-1], c)$
is a distinguished triangle. By TR3 there is a morphism of distinguished
triangles
$(\text{id}, \beta, \text{id}) : (A^\bullet, B^\bullet, C^\bullet, a, b, c) \to
(A^\bullet, W^\bullet[-1], C^\bullet, -i[-1], p[-1], c)$
which must be an isomorphism by Lemma \ref{lemma-third-isomorphism-triangle}.
This finishes the proof because
$0 \to A^\bullet \to W^\bullet[-1] \to C^\bullet \to 0$
is a termwise split short exact sequence of complexes
by the very construction of cones in Section \ref{section-cones}.
\end{proof}

\begin{remark}
\label{remark-homotopy-double}
Let $\mathcal{A}$ be an additive category with countable direct sums.
Let $\text{DoubleComp}(\mathcal{A})$ denote the category of double complexes
in $\mathcal{A}$, see
Homology, Section \ref{homology-section-double-complexes}.
We can use this category to construct two triangulated categories.
\begin{enumerate}
\item We can consider an object $A^{\bullet, \bullet}$ of
$\text{DoubleComp}(\mathcal{A})$ as a complex of complexes
as follows
$$
\ldots \to A^{\bullet, -1} \to A^{\bullet, 0} \to A^{\bullet, 1} \to \ldots
$$
and take the homotopy category $K_{first}(\text{DoubleComp}(\mathcal{A}))$
with the corresponding triangulated structure given by
Proposition \ref{proposition-homotopy-category-triangulated}.
By Homology, Remark
\ref{homology-remark-double-complex-complex-of-complexes-first} the functor
$$
\text{Tot} :
K_{first}(\text{DoubleComp}(\mathcal{A}))
\longrightarrow
K(\mathcal{A})
$$
is an exact functor of triangulated categories.
\item We can consider an object $A^{\bullet, \bullet}$ of
$\text{DoubleComp}(\mathcal{A})$ as a complex of complexes
as follows
$$
\ldots \to A^{-1, \bullet} \to A^{0, \bullet} \to A^{1, \bullet} \to \ldots
$$
and take the homotopy category $K_{second}(\text{DoubleComp}(\mathcal{A}))$
with the corresponding triangulated structure given by
Proposition \ref{proposition-homotopy-category-triangulated}.
By Homology, Remark
\ref{homology-remark-double-complex-complex-of-complexes-second} the functor
$$
\text{Tot} :
K_{second}(\text{DoubleComp}(\mathcal{A}))
\longrightarrow
K(\mathcal{A})
$$
is an exact functor of triangulated categories.
\end{enumerate}
\end{remark}

\begin{remark}
\label{remark-double-complex-as-tensor-product-of}
Let $\mathcal{A}$, $\mathcal{B}$, $\mathcal{C}$ be additive categories
and assume $\mathcal{C}$ has countable direct sums. Suppose that
$$
\otimes : \mathcal{A} \times \mathcal{B} \longrightarrow \mathcal{C},
\quad
(X, Y) \longmapsto X \otimes Y
$$
is a functor which is bilinear on morphisms. This determines a functor
$$
\text{Comp}(\mathcal{A}) \times \text{Comp}(\mathcal{B})
\longrightarrow
\text{DoubleComp}(\mathcal{C}), \quad
(X^\bullet, Y^\bullet)
\longmapsto
X^\bullet \otimes Y^\bullet
$$
See
Homology, Example \ref{homology-example-double-complex-as-tensor-product-of}.
\begin{enumerate}
\item For a fixed object $X^\bullet$ of $\text{Comp}(\mathcal{A})$
the functor
$$
K(\mathcal{B}) \longrightarrow K(\mathcal{C}), \quad
Y^\bullet \longmapsto \text{Tot}(X^\bullet \otimes Y^\bullet)
$$
is an exact functor of triangulated categories.
\item For a fixed object $Y^\bullet$ of $\text{Comp}(\mathcal{B})$
the functor
$$
K(\mathcal{A}) \longrightarrow K(\mathcal{C}), \quad
X^\bullet \longmapsto \text{Tot}(X^\bullet \otimes Y^\bullet)
$$
is an exact functor of triangulated categories.
\end{enumerate}
This follows from Remark \ref{remark-homotopy-double} since
the functors
$\text{Comp}(\mathcal{A}) \to \text{DoubleComp}(\mathcal{C})$,
$Y^\bullet \mapsto X^\bullet \otimes Y^\bullet$ and
$\text{Comp}(\mathcal{B}) \to \text{DoubleComp}(\mathcal{C})$,
$X^\bullet \mapsto X^\bullet \otimes Y^\bullet$
are immediately seen to be compatible with homotopies
and termwise split short exact sequences and hence induce
exact functors of triangulated categories
$$
K(\mathcal{B}) \to K_{first}(\text{DoubleComp}(\mathcal{C}))
\quad\text{and}\quad
K(\mathcal{A}) \to K_{second}(\text{DoubleComp}(\mathcal{C}))
$$
Observe that for the first of the two the isomorphism
$$
\text{Tot}(X^\bullet \otimes Y^\bullet[1]) \cong
\text{Tot}(X^\bullet \otimes Y^\bullet)[1]
$$
involves signs (this goes back to the signs chosen in
Homology, Remark \ref{homology-remark-shift-double-complex}).
\end{remark}






